\documentclass[../main.tex]{subfiles}
\begin{document}

\chapter{Cache Review}
\lessoninfo{
  video={Lesson 12, videos 1--45},
  textbook={H\&P \cite{hennessy2017} §2.1, §B.1--B.2},
  keywords={locality principle, temporal locality, spatial locality, cache hit, cache miss, AMAT, miss penalty, block offset, tag, valid bit, direct-mapped cache, set-associative cache, LRU, dirty bit},
}

Lessons 3--11 made processors complete ever more instructions per second,
yet every instruction has to be fetched from memory, many also load or
store data, and Lesson 1 showed that memory speed falls ever further behind
processor speed (the memory wall).  This lesson reviews the mechanism that
bridges the gap: the cache, a small and fast memory that serves most accesses
without waiting for the large and slow main memory, because programs access
memory with locality.  It covers how cache performance is measured, how a
cache is organized and searched, which block is evicted on a miss, and how
writes are handled; virtual memory (Lesson 13) and the advanced cache
techniques of Lesson 14 build directly on this material.

\section{The locality principle}
\label{sec:l12-locality}

Caches work because the memory accesses of a program are far from
random.\source{Lesson 12, videos 1--6; H\&P §2.1; \cite{denning2005}.}

\begin{definition}[Locality principle]
  The observation that what a program does in the near future is likely to be
  close to what it has done in the recent past is called the
  \term{locality principle}[局所性の原理].
\end{definition}

For memory references, \enquote{close} has two meanings.  Suppose a program
has just accessed address~$X$.

\begin{definition}[Temporal and spatial locality]
  A program exhibits \term{temporal locality}[時間的局所性] when it is likely
  to access the same address~$X$ again soon.  It exhibits
  \term{spatial locality}[空間的局所性] when it is likely to access addresses
  near~$X$ soon.
\end{definition}

Both kinds follow from the way programs are structured.  A loop executes the
same instructions over and over (temporal locality), and those instructions
are fetched from consecutive addresses (spatial locality).  A variable updated
in every iteration is accessed at the same address again and again; the
elements of an array processed in order lie at consecutive addresses.

\begin{example}
  \label{ex:l12-locality}
  The loop below copies R6 bytes from the buffer at address R4 to the
  buffer at address R5.
  \begin{lstlisting}[language={[x86masm]Assembler}, numbers=none]
Loop:  LB   R3, 0(R4)      ; R3 <- source byte
       SB   R3, 0(R5)      ; destination byte <- R3
       ADDI R4, R4, 1      ; next source byte
       ADDI R5, R5, 1      ; next destination byte
       ADDI R6, R6, -1     ; one byte less to copy
       BNE  R6, R0, Loop   ; repeat for all bytes
\end{lstlisting}
  The six instruction fetches have both kinds of locality: they go to
  consecutive addresses, and the same six addresses are fetched in every
  iteration.  The data accesses form two streams, one through each buffer.
  Every byte is read or written exactly once, so the data has no temporal
  locality.  It has spatial locality nevertheless: the accesses alternate
  between two buffers that may be far apart, but each of them is adjacent to
  the previous access to the same buffer.  Spatial locality only requires
  accesses near \emph{recently} accessed addresses, not near the one
  accessed last.\caution{Spatial locality is about distance, not
  regularity: a stream of accesses with a large constant stride is perfectly
  regular and has none.}
\end{example}

\section{Caches}
\label{sec:l12-caches}

Locality suggests how to give the processor fast access to a large
memory.\source{Lesson 12, videos 7--9; H\&P §B.1; \cite{wilkes1965}.}
Going to main memory for every fetch, load and store is slow.  Making all of
main memory as fast as the processor is not possible either: fast memory is
expensive per bit, and a larger memory is inherently slower to access.  The
remaining option is to keep a small amount of fast memory inside the
processor itself and to place in it the data that was accessed recently,
together with the data around it.  By temporal and spatial locality, most
accesses then find their data there.

\begin{definition}[Cache]
  A \term{cache}[キャッシュ] is a small, fast memory that holds copies of the
  recently accessed parts of a larger and slower memory, so that most
  accesses can be served from the copies.
\end{definition}

Every memory access of the processor first looks for its data in the cache
(\cref{fig:l12-cache-position}).  If the data is there, the access is a
\term{cache hit}[キャッシュヒット], and the data is read or written in the
cache.

Otherwise the access is a \term{cache miss}[キャッシュミス]: the data is
obtained from main memory and copied into the cache, so that further accesses
to it, and by spatial locality to its neighbours, will hit.  Because the
cache is consulted on every access, the lookup itself must be very fast, which
is a second reason, besides cost, why the cache has to be small.  A small
cache cannot hold everything, and misses are therefore
unavoidable.\sidenote{The idea is Wilkes's \enquote{slave memory} of 1965
\cite{wilkes1965}.  The first commercial machine to have one was the IBM
System/360 Model~85 of 1968, whose 16~KB buffer its designers named the
\emph{cache} \cite{liptay1968}.}

\begin{figure}[htbp]
  \centering
  % Processor, cache and main memory: every access looks in the cache first;
% only misses go on to main memory.
\begin{tikzpicture}[x=1cm,y=1cm, font=\footnotesize\sffamily, >={Stealth[length=1.8mm]},
  box/.style={draw, align=center, inner sep=3pt}]
  \node[box, fill=ln-box, minimum width=2.0cm, minimum height=1.2cm] (cpu) at (0,0)
    {\iflangja{プロセッサ}{processor}};
  \node[box, fill=ln-accent!15, minimum width=2.0cm, minimum height=1.0cm] (cache) at (4.0,0)
    {\iflangja{キャッシュ}{cache}\\[-1pt]{\scriptsize\iflangja{小容量・高速}{small, fast}}};
  \node[box, fill=white, minimum width=2.6cm, minimum height=1.8cm] (mem) at (8.2,0)
    {\iflangja{主記憶}{main memory}\\[-1pt]{\scriptsize\iflangja{大容量・低速}{large, slow}}};
  \draw[<->] (cpu) -- node[above, font=\sffamily\scriptsize] {\iflangja{すべてのアクセス}{every access}}
    node[below, font=\sffamily\scriptsize, text=ln-muted] {\iflangja{ヒット時間}{hit time}} (cache);
  \draw[<->, dashed] (cache) -- node[above, font=\sffamily\scriptsize] {\iflangja{ミス時のみ}{misses only}}
    node[below, font=\sffamily\scriptsize, text=ln-muted] {\iflangja{ミスペナルティ}{miss penalty}} (mem);
\end{tikzpicture}

  \caption{Every memory access goes to the cache; only misses go on to main
    memory.}
  \label{fig:l12-cache-position}
\end{figure}

\section{Cache performance}
\label{sec:l12-performance}

The benefit of a cache is measured by the average time a memory access
takes.\source{Lesson 12, videos 10--12; H\&P §B.2.}  The following
quantities enter it.

\begin{definition}[Hit rate, miss rate]
  The \term{hit rate}[ヒット率] is the fraction of memory accesses that are
  cache hits; the \term{miss rate}[ミス率] is the fraction that are misses.
  Hit rate $+$ miss rate $= 1$.
\end{definition}

\begin{definition}[Hit time, miss penalty]
  The \term{hit time}[ヒット時間] is the time needed to look up the cache and
  access the data in it.  The \term{miss penalty}[ミスペナルティ] is the
  \emph{additional} time a miss takes: fetching the data from main memory
  and placing it in the cache.
\end{definition}

A miss also pays the hit time, since the cache has to be searched before the
miss is detected; a miss therefore takes the hit time plus the miss penalty.
Weighting both cases by their frequency gives the
\term{average memory access time}[平均メモリアクセス時間]%
\index{AMAT|see{average memory access time}} (AMAT):
\begin{equation}
  \begin{aligned}
    \text{AMAT}
    &= (1-\text{miss rate}) \times \text{hit time}
       + \text{miss rate} \times (\text{hit time} + \text{miss penalty}) \\
    &= \text{hit time} + \text{miss rate} \times \text{miss penalty} .
  \end{aligned}
  \label{eq:l12-amat}
\end{equation}
\caution{The miss penalty is the \emph{extra} time of a miss: subtract the
hit time from the total time of a miss.}

A good AMAT requires three things.
\begin{itemize}
  \item A low hit time: the cache must be small and simple.
  \item A low miss rate: the cache must be large enough, and organized
    cleverly enough, to hold the data the program is about to use.
  \item A low miss penalty.  The miss penalty is dominated by the access time
    of main memory and is large; Lesson 14 reduces it with further levels of
    cache.
\end{itemize}
The first two goals conflict: a larger cache misses less often but takes
longer to search.  The best cache size balances them, and AMAT is the
measure by which they are balanced.

\begin{example}
  \label{ex:l12-amat}
  Main memory adds a miss penalty of
  \SI{60}{\nano\second}.\margin{At \SI{3}{\giga\hertz} this penalty is 180
  cycles; a 4-wide processor (Lesson 6) loses some 700 instruction slots
  while it waits.}  Three candidate
  caches are compared:
  \begin{center}\small
    \begin{tabular}{@{}lccc@{}}
      \toprule
      Cache & Hit time & Miss rate & AMAT \\
      \midrule
      Small  & \SI{1}{\nano\second}   & \SI{6}{\percent} & $1 + 0.06\times 60 = \SI{4.6}{\nano\second}$ \\
      Medium & \SI{1.5}{\nano\second} & \SI{3}{\percent} & $1.5 + 0.03\times 60 = \SI{3.3}{\nano\second}$ \\
      Large  & \SI{3}{\nano\second}   & \SI{2}{\percent} & $3 + 0.02\times 60 = \SI{4.2}{\nano\second}$ \\
      \bottomrule
    \end{tabular}
  \end{center}
  The medium cache is best.  The large cache misses less often, but its
  longer hit time costs more than the misses it saves.
  All three are an order of magnitude better than accessing main memory on
  every access.
\end{example}

\begin{note}
  A cache pays off only if a hit is much faster than an access to main
  memory and misses are rare.  With a hit time close to the access time of
  main memory, even a hit is about as slow as an access without a cache;
  with frequent misses, the miss penalty dominates.  In both cases AMAT
  approaches or even exceeds the access time of main memory alone.
\end{note}

\subsection{Caches in real processors}

Real processors have more than one level of cache.\source{Lesson 12, video
13; H\&P §2.1.}  The cache that is consulted first, and that directly serves
the reads and writes of the processor, is the smallest and fastest; only its
misses go on to further, larger cache levels (L2, L3; Lesson 14) and finally
to main memory.  This first level is called the \term{L1 cache}[L1
キャッシュ].  L1 caches are typically 16--64~KB (1~KB $= 2^{10}$ bytes):
large enough that about \SI{90}{\percent} of accesses hit, and small enough
that a hit takes only one to three clock cycles.\margin{Sizes today:
Skylake has a 32~KB L1 per core, a 256~KB L2 and an 8~MB shared L3; Zen~3 has
32~KB, 512~KB and 32~MB (7-cpu.com).}

\section{Blocks and lines}
\label{sec:l12-blocks}

To determine quickly whether data is in the cache, the cache is organized as
a table of fixed-size entries.\source{Lesson 12, videos 14--18; H\&P §B.1.}
Data is moved between main memory and the cache not byte by byte but in
units called blocks.

\begin{definition}[Cache block]
  A \term{cache block}[キャッシュブロック] is a run of consecutive bytes of
  memory that is copied into the cache, and evicted from it, as a unit.  The
  number of bytes in a block is the \term{block size}[ブロックサイズ].
\end{definition}

The block size trades spatial locality against cache space.
\begin{itemize}
  \item If blocks are very small (in the extreme, one byte), an access to
    the next byte, which spatial locality makes likely, misses again, and
    each byte needs its own bookkeeping.
  \item If blocks are very large, most of a block is data that will not be
    used soon.  It occupies space in which other recently used data could
    have been kept, and fewer blocks fit into the cache.
\end{itemize}
Typical block sizes are 32--128 bytes.\margin{Real block sizes: 64 bytes on
every x86-64 processor and on most Arm cores, 128 bytes on Apple's M1 and
M2.}  The number of blocks a
cache can hold is its size divided by the block size: a
16~KB cache with 64-byte blocks holds
$2^{14}/2^{6} = 256$ blocks.

\subsection{Where blocks start}

Blocks do not start at arbitrary addresses.  If they did, a block covering
addresses 100--227 and another covering 200--327 could both be in the cache,
an address could lie in several cached blocks, and finding out whether an
address is cached would require comparing address ranges.  Instead, memory
is divided into fixed, non-overlapping blocks that start at multiples of the
block size: with 128-byte blocks, at addresses 0, 128, 256, \ldots\
(\cref{fig:l12-blocks}).  Every address then belongs to exactly one block.

\begin{definition}[Cache line]
  A \term{cache line}[キャッシュライン] is a slot in the cache that can hold
  one block.  Memory is divided into blocks and the cache into lines; a line
  has the size of a block, and holds different blocks at different times.
\end{definition}

Block sizes, and therefore line sizes, are always powers of
two.\caution{The lecture keeps \emph{block} (in memory) and \emph{line} (its
slot in the cache) apart; H\&P and vendor manuals call both a cache line.}
Only then
can the block of an address and the position within the block be found by
selecting bits of the address, as the next section shows, rather than by a
division.

\begin{figure}[htbp]
  \centering
  % Memory is divided into 128-byte blocks starting at multiples of the block
% size; the cache into lines holding one block each, with valid bit and tag.
\begin{tikzpicture}[x=1cm,y=1cm, font=\footnotesize\sffamily, >={Stealth[length=1.8mm]}]
  % --- address fields
  \draw[fill=ln-accent!15] (1.2,0) rectangle (6.4,0.55);
  \node at (3.8,0.275) {\iflangja{ブロック番号}{block number}\ (31\,\dots\,7)};
  \draw[fill=ln-box] (6.4,0) rectangle (9.6,0.55);
  \node at (8.0,0.275) {\iflangja{ブロックオフセット}{block offset}\ (6\,\dots\,0)};
  \node[anchor=east] at (1.15,0.275) {\iflangja{アドレス}{address}};
  % --- main memory
  \node[font=\sffamily\scriptsize, text=ln-muted, anchor=south] at (2.3,-1.0) {\iflangja{主記憶}{main memory}};
  \foreach \i/\a in {0/0x000,1/0x080,2/0x100,3/0x180,4/0x200} {
    \draw[fill=white] (1.2,-1.0-0.5*\i) rectangle (3.4,-1.5-0.5*\i);
    \node[font=\sffamily\scriptsize\ttfamily, anchor=east, text=ln-muted] at (1.15,-1.0-0.5*\i) {\a};
    \node at (2.3,-1.25-0.5*\i) {\iflangja{ブロック}{block}\ \i};
  }
  \fill[ln-accent!15] (1.2,-1.5) rectangle (3.4,-2.0);
  \fill[ln-accent!15] (1.2,-3.0) rectangle (3.4,-3.5);
  \foreach \i in {1,4} {
    \draw (1.2,-1.0-0.5*\i) rectangle (3.4,-1.5-0.5*\i);
    \node at (2.3,-1.25-0.5*\i) {\iflangja{ブロック}{block}\ \i};
  }
  \node at (2.3,-3.75) {$\vdots$};
  % --- cache lines
  \node[font=\sffamily\scriptsize, text=ln-muted, anchor=south] at (6.35,-1.0) {V};
  \node[font=\sffamily\scriptsize, text=ln-muted, anchor=south] at (7.15,-1.0) {\iflangja{タグ}{tag}};
  \node[font=\sffamily\scriptsize, text=ln-muted, anchor=south] at (8.9,-1.0) {\iflangja{データ（128 バイト）}{data (128 bytes)}};
  \foreach \i/\v/\t/\d in {0/1/4/{\iflangja{ブロック 4 の写し}{copy of block 4}},
                           1/0/?/?,
                           2/1/1/{\iflangja{ブロック 1 の写し}{copy of block 1}},
                           3/0/?/?} {
    \draw[fill=white] (6.1,-1.0-0.5*\i) rectangle (6.6,-1.5-0.5*\i);
    \draw[fill=white] (6.6,-1.0-0.5*\i) rectangle (7.7,-1.5-0.5*\i);
    \draw[fill=ln-box] (7.7,-1.0-0.5*\i) rectangle (10.1,-1.5-0.5*\i);
    \node at (6.35,-1.25-0.5*\i) {\v};
    \node at (7.15,-1.25-0.5*\i) {\t};
    \node[font=\sffamily\scriptsize] at (8.9,-1.25-0.5*\i) {\d};
    \node[font=\sffamily\scriptsize, text=ln-muted, anchor=west] at (10.15,-1.25-0.5*\i) {\iflangja{ライン}{line}\ \i};
  }
  \node[font=\sffamily\scriptsize, text=ln-muted, anchor=north] at (8.1,-3.05) {\iflangja{キャッシュ}{cache}};
  % --- copies
  \draw[->, dashed, ln-accent] (3.4,-1.75) .. controls (4.8,-1.75) and (4.6,-2.25) .. (6.1,-2.25);
  \draw[->, dashed, ln-accent] (3.4,-3.25) .. controls (4.9,-3.25) and (4.4,-1.25) .. (6.1,-1.25);
\end{tikzpicture}

  \caption{Main memory divided into 128-byte blocks and a cache of
    four lines.  Lines~0 and~2 hold copies of blocks~4 and~1; lines~1 and~3
    hold nothing valid.}
  \label{fig:l12-blocks}
\end{figure}

\section{Block number, tag and valid bit}
\label{sec:l12-tags}

With a block size of $2^{k}$ bytes, an address~$A$ splits into two
fields.\source{Lesson 12, videos 19--23; H\&P §B.1.}

\begin{definition}[Block offset and block number]
  The low-order $k$ bits of the address are its
  \term{block offset}[ブロックオフセット], the position of the byte within
  its block.  The remaining high-order bits are its
  \term{block number}[ブロック番号], which identifies the block:
  \begin{equation}
    \text{block number} = \bigl\lfloor A / 2^{k} \bigr\rfloor ,
    \qquad
    \text{block offset} = A \bmod 2^{k} .
    \label{eq:l12-block-number}
  \end{equation}
\end{definition}

\begin{example}
  \label{ex:l12-block-number}
  With 128-byte blocks ($k=7$), the address \texttt{0x0040A7C6} has
  block offset $\texttt{0xC6} \mathbin{\&} \texttt{0x7F} = \texttt{0x46} = 70$
  and block number $\texttt{0x0040A7C6} \gg 7 = \texttt{0x814F}$.  The byte
  is at position~70 of block \texttt{0x814F}, which covers the addresses
  \texttt{0x0040A780} to \texttt{0x0040A7FF}.
\end{example}

\subsection{Tags}

Since a line holds different blocks at different times, each line must
record which block it currently holds.  This record is the
\term{tag}[タグ] of the line: the block number of the block in the line, or,
as the following sections show, the part of the block number that is not
already implied by the position of the line.  A lookup computes the block
number of the address and compares it with the tags of the lines in which
the block could be.  A match is a hit, and the block offset then selects the
byte within the line.

\subsection{The valid bit}

When the processor is turned on, the lines of the cache contain arbitrary
bits, including arbitrary tags.  A tag that happens to match a block number
would produce a hit on data that was never loaded.

\begin{definition}[Valid bit]
  The \term{valid bit}[有効ビット] of a line indicates whether the line holds
  a block.  All valid bits are cleared at start-up, and a line's valid bit is
  set when a block is loaded into it.
\end{definition}

Emptying a line, or the whole cache, thus only requires clearing valid bits;
the tags and data can be left as they are.  The one exception is a block
that has been written in the cache but not yet in main memory
(\cref{sec:l12-write}): it must be written back first.\margin{x86 has both
operations: \texttt{WBINVD} writes dirty blocks back before clearing the
valid bits, \texttt{INVD} clears them and discards the data.}

A lookup is a hit only if the line is valid and its stored tag matches the
tag field of the address:
\begin{equation}
  \text{hit} \;=\; V \wedge (\text{stored tag} = \text{tag of the address}) .
  \label{eq:l12-hit}
\end{equation}
As long as no part of the block number selects the line, as in this section,
the tag field is the whole block number; the organizations of the following
sections use some of its bits to select the line and keep only the remaining
bits in the tag.

\begin{keypoint}
  An address is a block number and an offset.  The tag identifies which
  block a line holds, the valid bit whether it holds one at all, and the
  offset picks the byte.  Where a block may be placed, the subject of the
  next section, decides \emph{which} lines have to be searched.
\end{keypoint}

\section{Where a block can be placed}
\label{sec:l12-placement}

Caches differ in how many of their lines a given block may be placed
in.\source{Lesson 12, video 24; H\&P §B.1.}  This number is the
\term{associativity}[連想度] of the cache.  It determines how many tags a
lookup must compare, and how often different blocks compete for the same
lines.  Three organizations are distinguished (\cref{fig:l12-placement}):
\begin{description}
  \item[Direct-mapped] each block can be placed in exactly one line
    (\cref{sec:l12-dm}).
  \item[Set-associative] each block can be placed in any of $N$ lines that
    form a set (\cref{sec:l12-sa}).
  \item[Fully associative] each block can be placed in any line of the
    cache (\cref{sec:l12-fa}).\margin{Typical associativities: Skylake L1
    8-way, L2 4-way, L3 16-way; Zen~3 L1 8-way, L2 8-way, L3 16-way.}
\end{description}

\begin{figure}[htbp]
  \centering
  % Where block number 13 may be placed in an 8-line cache that is
% direct-mapped, 2-way set-associative, or fully associative.
\begin{tikzpicture}[x=1cm,y=1cm, font=\footnotesize\sffamily]
  \def\rh{0.34}
  \foreach \cx/\ttl/\how in {
      1.7/{\iflangja{ダイレクトマップ}{direct-mapped}}/{\iflangja{ライン}{line} $13 \bmod 8 = 5$},
      5.5/{\iflangja{2 ウェイセットアソシアティブ}{2-way set-associative}}/{\iflangja{セット}{set} $13 \bmod 4 = 1$},
      9.3/{\iflangja{フルアソシアティブ}{fully associative}}/{\iflangja{任意のライン}{any line}}} {
    \node[font=\sffamily\scriptsize\bfseries, anchor=south, text width=3.4cm, align=center] at (\cx,0.08) {\ttl};
    \node[font=\sffamily\scriptsize, anchor=north] at (\cx,-8*\rh-0.1) {\how};
    \foreach \i in {0,...,7} {
      \node[font=\sffamily\scriptsize, text=ln-muted, anchor=east] at (\cx-0.8,-\rh*\i-\rh/2) {\i};
    }
  }
  % highlighted lines
  \fill[ln-accent!15] (0.9,-5*\rh) rectangle (2.5,-6*\rh);
  \fill[ln-accent!15] (4.7,-2*\rh) rectangle (6.3,-4*\rh);
  \fill[ln-accent!15] (8.5,0) rectangle (10.1,-8*\rh);
  % grids
  \foreach \cx in {1.7,5.5,9.3} {
    \foreach \i in {0,...,7} { \draw (\cx-0.8,-\rh*\i) rectangle (\cx+0.8,-\rh*\i-\rh); }
  }
  % set boundaries and labels in the 2-way cache
  \foreach \s in {0,...,3} {
    \draw[thick] (4.7,-2*\s*\rh) rectangle (6.3,-2*\s*\rh-2*\rh);
    \node[font=\sffamily\scriptsize, text=ln-muted, anchor=west] at (6.35,-2*\s*\rh-\rh) {\iflangja{セット}{set}\ \s};
  }
\end{tikzpicture}

  \caption{The lines in which block number~13 may be placed in an 8-line
    cache of each organization.}
  \label{fig:l12-placement}
\end{figure}

\section{Direct-mapped caches}
\label{sec:l12-dm}

The simplest of the three organizations gives every block a single line of
the cache.\source{Lesson 12, videos 25--28; H\&P §B.1.}  In a
\term{direct-mapped cache}[ダイレクトマップキャッシュ] every block has exactly
one line in which it can be placed, and that line is selected by the
low-order bits of the block number.

\begin{definition}[Index]
  The \term{index}[インデックス] is the group of address bits that selects
  where in the cache a block may be placed.  In a direct-mapped cache with
  $2^{n}$ lines it consists of the $n$ low-order bits of the block number.
\end{definition}

The address is thus split into three fields (\cref{fig:l12-direct-mapped}):
the offset in the $k$ lowest bits, above it the index in the next $n$ bits,
and the tag in all remaining bits.  The line stores only these tag bits, not
the whole block number: all blocks that can be placed in line~$i$ have
index~$i$, so the index part of their block numbers is known from the line
itself.

A lookup reads the one line selected by the index, compares its tag with the
tag of the address, and checks the valid bit (\cref{eq:l12-hit}).

\begin{figure}[htbp]
  \centering
  % Direct-mapped lookup for a 64 KB cache with 128-byte blocks and 32-bit
% addresses: the index selects one line, one comparator checks the tag.
\begin{tikzpicture}[x=1cm,y=1cm, font=\footnotesize\sffamily, >={Stealth[length=1.8mm]}]
  % --- address fields
  %% The Japanese index label filled its box edge to edge, so the Japanese
  %% tag/index boundary sits further left (the index arrow at x=3.3 and the
  %% tag arrow at x=1.5 stay inside their boxes).
  \iflangja{%
    \draw[fill=ln-accent!15] (0,0) rectangle (2.6,0.55);
    \node at (1.3,0.275) {タグ\ (31\,\dots\,16)};
    \draw[fill=ln-box] (2.6,0) rectangle (5.6,0.55);
    \node at (4.1,0.275) {インデックス\ (15\,\dots\,7)};
  }{%
    \draw[fill=ln-accent!15] (0,0) rectangle (3.0,0.55);
    \node at (1.5,0.275) {tag\ (31\,\dots\,16)};
    \draw[fill=ln-box] (3.0,0) rectangle (5.6,0.55);
    \node at (4.3,0.275) {index\ (15\,\dots\,7)};
  }
  \draw[fill=white] (5.6,0) rectangle (8.2,0.55);
  \node at (6.9,0.275) {\iflangja{オフセット}{offset}\ (6\,\dots\,0)};
  % --- table: V | tag | data, 6 drawn rows
  \node[font=\sffamily\scriptsize, text=ln-muted, anchor=south] at (3.85,-1.0) {V};
  \node[font=\sffamily\scriptsize, text=ln-muted, anchor=south] at (4.9,-1.0) {\iflangja{タグ}{tag}};
  \node[font=\sffamily\scriptsize, text=ln-muted, anchor=south] at (7.2,-1.0) {\iflangja{データ}{data}};
  \fill[ln-accent!15] (3.6,-2.35) rectangle (8.7,-2.8);
  \foreach \i/\lab in {0/0,1/1,2/{},3/335,4/{},5/511} {
    \draw (3.6,-1.0-0.45*\i) rectangle (4.1,-1.45-0.45*\i);
    \draw (4.1,-1.0-0.45*\i) rectangle (5.7,-1.45-0.45*\i);
    \draw (5.7,-1.0-0.45*\i) rectangle (8.7,-1.45-0.45*\i);
    \node[font=\sffamily\scriptsize, text=ln-muted, anchor=west] at (8.75,-1.225-0.45*\i) {\lab};
  }
  \foreach \i in {2,4} {
    \node[font=\small, text=ln-muted, inner sep=0pt] at (4.9,-1.225-0.45*\i) {$\vdots$};
    \node[font=\small, text=ln-muted, inner sep=0pt] at (7.2,-1.225-0.45*\i) {$\vdots$};
  }
  % --- index selects a line
  \draw[->] (3.3,0) -- (3.3,-2.575) -- (3.6,-2.575);
  % --- tag comparison and valid bit
  \node[draw, circle, fill=white, inner sep=1pt, minimum size=0.5cm] (eq) at (4.9,-4.25) {$=$};
  \node[draw, fill=ln-box, minimum width=0.7cm, minimum height=0.5cm, inner sep=1pt, font=\sffamily\scriptsize] (and) at (3.85,-4.25) {AND};
  \draw[->] (4.9,-3.7) -- (eq.north);
  \draw[->] (3.85,-3.7) -- (and.north);
  \draw[->] (1.5,0) -- (1.5,-5.0) -- (4.9,-5.0) -- (eq.south);
  \draw[->] (eq.west) -- (and.east);
  \draw[->] (and.west) -- (2.6,-4.25) node[left] {\iflangja{ヒット}{hit}};
  % --- data selection by offset
  \node[draw, fill=white, minimum width=2.0cm, minimum height=0.5cm, inner sep=1pt, font=\sffamily\scriptsize] (sel) at (7.2,-4.25)
    {\iflangja{バイトを選択}{select byte}};
  \draw[->] (7.2,-3.7) -- (sel.north);
  \draw[->] (8.2,0.275) -- (9.6,0.275) -- (9.6,-4.25) -- (sel.east);
  \draw[->] (sel.south) -- (7.2,-5.0) node[below] {\iflangja{データ}{data}};
\end{tikzpicture}

  \caption{Lookup in a direct-mapped 64~KB cache with
    128-byte blocks.  The index selects one of 512 lines; a single
    comparator checks its tag.}
  \label{fig:l12-direct-mapped}
\end{figure}

\begin{example}
  \label{ex:l12-dm-fields}
  For a direct-mapped 64~KB cache with 128-byte blocks
  and 32-bit addresses there are $512 = 2^{9}$ lines: offset bits 6--0,
  index bits 15--7, tag bits 31--16.  The address \texttt{0x0040A7C6} of
  \cref{ex:l12-block-number}, block number \texttt{0x814F}, has index
  $\texttt{0x814F} \mathbin{\&} \texttt{0x1FF} = \texttt{0x14F} = 335$ and
  tag $\texttt{0x814F} \gg 9 = \texttt{0x0040}$.  It hits if line~335 is
  valid and holds tag \texttt{0x0040}.
\end{example}

\subsection{Advantages and disadvantages}

The strength of a direct-mapped cache is that only one line has to be
examined.  A single comparator suffices, so the lookup is fast, cheap and
energy-efficient: the hit time is as low as it can be for a cache of a given
size.

Its weakness is that blocks with the same index compete for a single line.
If a program alternates between two such blocks, each access evicts the block
the other access needs, and every access misses, even when the rest of the
cache is empty.  These conflicts raise the miss rate.\margin{Misses caused
by this competition are \emph{conflict misses}, one of the three causes of
misses classified in Lesson 14.}

\begin{example}
  \label{ex:l12-dm-conflict}
  An 8-line direct-mapped cache with 16-byte blocks uses offset bits
  3--0, index bits 6--4 and tag bits 31--7.  A program alternately accesses
  addresses \texttt{0x1000} (block \texttt{0x100}) and \texttt{0x1084}
  (block \texttt{0x108}).  Both blocks have index~0, with tags \texttt{0x20}
  and \texttt{0x21}.  Each access replaces the block loaded by the previous
  one, so every access is a miss, however long the program continues.
\end{example}

\section{Set-associative caches}
\label{sec:l12-sa}

A set-associative organization relaxes the one-line
restriction.\source{Lesson 12, videos 29--31; H\&P §B.1.}

\begin{definition}[Set-associative cache]
  A \term{set-associative cache}[セットアソシアティブキャッシュ] divides its
  lines into \emph{sets} of $N$ lines each; a block can be placed in any
  line of one particular set.  Such a cache is called $N$-way
  set-associative.
\end{definition}

The index now selects a set rather than a line.  A cache with $2^{s}$ sets
has an $s$-bit index, and the tag again consists of all bits above the
index.  A lookup compares the tag of the address with the tags of all $N$
lines of the selected set, in parallel, and hits if any valid line matches
(\cref{fig:l12-set-associative}); the matching line supplies the data.

\begin{figure}[htbp]
  \centering
  % Lookup in a 2-way set-associative cache with 4 sets (8 lines, 16-byte
% blocks, 32-bit addresses): the index selects a set, both tags are compared,
% and each comparison is ANDed with the valid bit of its line.
\begin{tikzpicture}[x=1cm,y=1cm, font=\footnotesize\sffamily, >={Stealth[length=1.8mm]},
  gate/.style={draw, fill=ln-box, minimum width=0.7cm, minimum height=0.5cm, inner sep=1pt, font=\sffamily\scriptsize},
  cmp/.style={draw, circle, fill=white, inner sep=1pt, minimum size=0.5cm}]
  % --- address fields
  %% The Japanese index label filled its box edge to edge, so the Japanese
  %% tag/index boundary sits further left (the index line at x=5.4 and the
  %% tag line at x=1.0 stay inside their boxes).
  \iflangja{%
    \draw[fill=ln-accent!15] (0.6,0) rectangle (3.6,0.55);
    \node at (2.1,0.275) {タグ\ (31\,\dots\,6)};
    \draw[fill=ln-box] (3.6,0) rectangle (6.4,0.55);
    \node at (5.0,0.275) {インデックス\ (5\,\dots\,4)};
  }{%
    \draw[fill=ln-accent!15] (0.6,0) rectangle (4.0,0.55);
    \node at (2.3,0.275) {tag\ (31\,\dots\,6)};
    \draw[fill=ln-box] (4.0,0) rectangle (6.4,0.55);
    \node at (5.2,0.275) {index\ (5\,\dots\,4)};
  }
  \draw[fill=white] (6.4,0) rectangle (8.8,0.55);
  \node at (7.6,0.275) {\iflangja{オフセット}{offset}\ (3\,\dots\,0)};
  % --- two ways, four sets; set 0 selected
  \foreach \x in {1.6,5.8} {
    \fill[ln-accent!15] (\x,-1.1) rectangle (\x+3.4,-1.55);
    \node[font=\sffamily\scriptsize, text=ln-muted, anchor=south] at (\x+0.2,-1.1) {V};
    \node[font=\sffamily\scriptsize, text=ln-muted, anchor=south] at (\x+1.1,-1.1) {\iflangja{タグ}{tag}};
    \node[font=\sffamily\scriptsize, text=ln-muted, anchor=south] at (\x+2.6,-1.1) {\iflangja{データ}{data}};
    \foreach \i in {0,...,3} {
      \draw (\x,-1.1-0.45*\i) rectangle (\x+0.4,-1.55-0.45*\i);
      \draw (\x+0.4,-1.1-0.45*\i) rectangle (\x+1.8,-1.55-0.45*\i);
      \draw (\x+1.8,-1.1-0.45*\i) rectangle (\x+3.4,-1.55-0.45*\i);
    }
  }
  \node[font=\sffamily\scriptsize, text=ln-muted, anchor=north] at (4.2,-2.95) {\iflangja{ウェイ 0}{way 0}};
  \node[font=\sffamily\scriptsize, text=ln-muted, anchor=north] at (8.4,-2.95) {\iflangja{ウェイ 1}{way 1}};
  \foreach \i in {0,...,3} {
    \node[font=\sffamily\scriptsize, text=ln-muted, anchor=west] at (9.25,-1.325-0.45*\i) {\iflangja{セット}{set}\ \i};
  }
  % --- index selects a set in both ways
  \draw (5.4,0) -- (5.4,-1.325);
  \fill (5.4,-1.325) circle (1.2pt);
  \draw[->] (5.4,-1.325) -- (5.0,-1.325);
  \draw[->] (5.4,-1.325) -- (5.8,-1.325);
  % --- per way: tag comparator, AND with the valid bit
  \foreach \x/\w in {1.6/0,5.8/1} {
    \node[cmp] (c\w) at (\x+1.2,-3.6) {$=$};
    \node[gate] (a\w) at (\x+0.4,-4.35) {AND};
    \draw[->] (\x+1.2,-2.9) -- (c\w.north);
    \draw[->] (\x+0.2,-2.9) -- (\x+0.2,-4.1);
    \draw[->] (c\w.west) -| (\x+0.6,-4.1);
  }
  % --- tag of the address to both comparators
  \draw (1.0,0) -- (1.0,-5.5) -- (7.0,-5.5);
  \draw[->] (7.0,-5.5) -- (c1.south);
  \fill (2.8,-5.5) circle (1.2pt);
  \draw[->] (2.8,-5.5) -- (c0.south);
  % --- either way hits
  \node[gate] (or) at (4.5,-4.35) {OR};
  \draw[->] (a0.east) -- (or.west);
  \draw[->] (a1.west) -- (or.east);
  \draw[->] (or.south) -- (4.5,-4.95) node[below] {\iflangja{ヒット}{hit}};
\end{tikzpicture}

  \caption{Lookup in a 2-way set-associative cache with four sets, the
    organization of \cref{ex:l12-sa-conflict}.  The index (here~0, as in
    that example) selects a set; each of its two lines hits if it is valid
    and its tag matches.}
  \label{fig:l12-set-associative}
\end{figure}

\begin{example}
  \label{ex:l12-sa-conflict}
  Organize the 8 lines of \cref{ex:l12-dm-conflict} as a 2-way
  set-associative cache.  It has 4 sets, so the index is bits 5--4 and the
  tag bits 31--6.  Block \texttt{0x100} goes to set~0 with tag
  \texttt{0x40}, and block \texttt{0x108} to set~0 with tag \texttt{0x42}.
  The two blocks still share a set, but the set has two lines: after the
  first access to each block has missed and loaded it, every further access
  hits.
\end{example}

Compared with a direct-mapped cache of the same size, an $N$-way
set-associative cache needs $N$ comparators and a selection of the matching
line's data, which makes the hit time longer, but it suffers fewer conflicts
and so has a lower miss rate.\margin{Why 32~KB 8-way L1 caches are common: 64
sets of 64-byte blocks span 4~KB, so index and offset fit in the page offset
and the lookup can start before translation (Lesson 14).}

\section{Fully associative caches}
\label{sec:l12-fa}

At the other extreme, a block may be placed
anywhere.\source{Lesson 12, videos 32--33; H\&P §B.1.}

\begin{definition}[Fully associative cache]
  In a \term{fully associative cache}[フルアソシアティブキャッシュ] any
  block can be placed in any line.
\end{definition}

There is nothing to select, so there is no index: the tag is the entire
block number, and a lookup compares it with the tag of every line in the
cache.  Two blocks never have to compete for a line while another line is
free, but a comparator is needed for every line, which makes a large fully
associative cache slow and expensive.\margin{Fully associative structures
stay small: TLBs of 64--512 entries (Lesson 13) are fully associative or
nearly so; caches are not.}

\subsection{One scheme for all three}

Direct-mapped and fully associative caches are the two extremes of
set-associative caches.  A direct-mapped cache is 1-way set-associative, with
as many sets as lines; a fully associative cache has a single set containing
all lines.  For every organization, with addresses of $a$ bits,
\begin{equation}
  \begin{aligned}
    \text{offset bits} &= \log_2(\text{block size}), \\
    \text{index bits}  &= \log_2(\text{number of sets})
                        = \log_2\bigl(\text{lines}/N\bigr), \\
    \text{tag bits}    &= a - \text{index bits} - \text{offset bits} ,
  \end{aligned}
  \label{eq:l12-fields}
\end{equation}
where $N=1$ for a direct-mapped cache (index bits $= \log_2(\text{lines})$)
and $N = \text{lines}$ for a fully associative cache (0 index bits).  For
the 64~KB cache with 128-byte blocks of
\cref{ex:l12-dm-fields}, the tag grows from 16 bits when direct-mapped to
25 bits when fully associative; \cref{tab:l12-organizations} in
\cref{sec:l12-summary} lists the intermediate cases.\tip{A cache of $C$ bytes has
$a - \log_2(C/N)$ tag bits.  Each doubling of $N$ moves one bit from the index
to the tag.}

\begin{keypoint}
  Associativity trades hit time against miss rate.  More ways per set means
  more tag comparisons on every access, but fewer conflicts, because a block
  is evicted only when every line of its set is occupied by another block.
\end{keypoint}

\section{Cache replacement}
\label{sec:l12-replacement}

On a miss, the new block must be placed in one of the lines allowed for
it.\source{Lesson 12, videos 34--36; H\&P §B.1.}  If one of those lines is
not valid, that line is used.  If all of them hold valid blocks, one block
has to be evicted.  A direct-mapped cache has no choice; in set-associative
and fully associative caches the victim is chosen by the
\term{replacement policy}[置換アルゴリズム].
\begin{description}
  \item[Random] evict a randomly chosen block of the set.  Simple, but it
    ignores how the blocks are being used.
  \item[FIFO] (first in, first out) evict the block that has been in the
    cache longest, regardless of whether it is still in use.
  \item[LRU] evict the block whose most recent access lies furthest in the
    past.
\end{description}

\begin{definition}[Least recently used]
  The \term{least recently used}[LRU]\index{LRU|see{least recently used}}
  (LRU) policy evicts, among the blocks that may be replaced, the one that
  has not been accessed for the longest time.
\end{definition}

LRU exploits temporal locality: a block accessed recently is likely to be
accessed again soon, so the least recently used block is the one least likely
to be needed.

\subsection{Implementing LRU with counters}

LRU can be implemented with a counter per line.  In an $N$-way
set-associative cache each line has a $\log_2 N$-bit counter.  At start-up
the counters of each set are initialized to the values $0, 1, \ldots, N-1$,
in any order, and the update rule below keeps it that way: every value
occurs exactly once in each set.\tip{The counters of a set stay a permutation
of $0,\ldots,N-1$: a value that appears twice in a row of a trace means an
update was missed.}  The line with counter~0 is the least
recently used and the one with counter $N-1$ the most recently used.  On
every access to a line whose counter is~$c$:
\begin{enumerate}
  \item every counter in the set that is greater than $c$ is decremented;
  \item the counter of the accessed line is set to $N-1$.
\end{enumerate}
On a miss, the victim is the line with counter~0; once the new block is in
it, the access is treated like any other and the counters are updated as
above.  \Cref{tab:l12-lru-trace} traces one 4-way set.\caution{Update the
counters on \emph{every} access, hits included; updating them only on loads
gives FIFO.}

\begin{table}[htbp]
  \centering\small
  \caption{LRU counters of one 4-way set holding blocks A--D, shown as
    block:counter for each of its four lines.}
  \label{tab:l12-lru-trace}
  \begin{tabular}{@{}llcccc@{}}
    \toprule
    Access & Result & Line 0 & Line 1 & Line 2 & Line 3 \\
    \midrule
    (initial) &                     & A:0 & B:1 & C:2 & D:3 \\
    C         & hit                 & A:0 & B:1 & C:3 & D:2 \\
    A         & hit                 & A:3 & B:0 & C:2 & D:1 \\
    E         & miss, B evicted     & A:2 & E:3 & C:1 & D:0 \\
    D         & hit                 & A:1 & E:2 & C:0 & D:3 \\
    B         & miss, C evicted     & A:0 & E:1 & B:3 & D:2 \\
    \bottomrule
  \end{tabular}
\end{table}

The cost of this implementation is $\log_2 N$ bits per line and, more
importantly, activity on every access: even a hit may change all $N$
counters of the set.  For highly associative caches this costs considerable
hardware and energy, which motivates the cheaper approximations of LRU in
Lesson 14.

\section{Write policy}
\label{sec:l12-write}

Writes raise two questions that reads do not.\source{Lesson 12, videos
37--40; H\&P §B.1.}

\begin{enumerate}
  \item \emph{Is a block brought into the cache on a write miss?}  A cache
    that brings it in follows the
    \term{write-allocate policy}[ライトアロケート方式]; a no-write-allocate
    cache writes the value to main memory only.  Most caches are
    write-allocate: by locality, a block that has just been written is likely
    to be read or written again soon.
  \item \emph{When does main memory receive the written value?}  Here two
    policies exist.
\end{enumerate}

\begin{definition}[Write-through and write-back]
  Under the \term{write-through policy}[ライトスルー方式] a write updates
  the block in the cache and immediately also the block in main memory.

  Under the \term{write-back policy}[ライトバック方式] a write updates only
  the block in the cache; main memory is updated when the block is evicted
  from the cache.
\end{definition}

Write-through keeps main memory always up to date, but every write then pays
the access time of main memory, which is what the cache was meant to avoid.
Caches therefore use write-back: many writes to a block while it is in the
cache cost at most a single write to main memory, at eviction.  A write-back
cache is normally also write-allocate; otherwise every write to a block that
is not in the cache would go to main memory, however often the block is
written.\caution{The two choices are independent: write-through or write-back
says \emph{when} memory is updated, write-allocate or not \emph{what a write
miss does}.}

Writing every evicted block back would, however, also write blocks that were
only read and are identical to their copies in memory.

\begin{definition}[Dirty bit]
  The \term{dirty bit}[ダーティビット] of a line indicates whether its block
  has been written since it was loaded.  It is cleared when a block is loaded
  and set when the block is written.  On eviction, a dirty block is written to
  main memory; a clean block is simply discarded.
\end{definition}

\begin{example}
  \label{ex:l12-write-back}
  Blocks X and Y map to the same line of a direct-mapped cache, which is
  initially invalid.  \Cref{tab:l12-write-back} follows a sequence of
  accesses under write-back.  The two consecutive writes to Y cost a single
  memory write, at step~6, and evicting Y at step~8 costs none, because Y
  was not written after it was loaded at step~7.  Write-back needs 3 memory
  writes in total; write-through would need one per write,
  i.e.~4.\sidenote{Write-through is not extinct: the Pentium~4 and AMD's
  Bulldozer wrote their L1 data caches through to the next level and absorbed
  the writes in a buffer.  Such an L1 needs no dirty bits and never writes
  back; both families later returned to write-back L1 caches.}
\end{example}

\begin{table}[htbp]
  \centering\small
  \caption{A single line under the write-back policy
    (\cref{ex:l12-write-back}).  D is the dirty bit after the access.}
  \label{tab:l12-write-back}
  \begin{tabular}{@{}cllll@{}}
    \toprule
    Step & Access  & Result & Memory operations        & Line after \\
    \midrule
    1    & read X  & miss   & read X                    & X, D$=$0 \\
    2    & write X & hit    & ---                       & X, D$=$1 \\
    3    & read Y  & miss   & write X back, read Y      & Y, D$=$0 \\
    4    & write Y & hit    & ---                       & Y, D$=$1 \\
    5    & write Y & hit    & ---                       & Y, D$=$1 \\
    6    & write X & miss   & write Y back, read X      & X, D$=$1 \\
    7    & read Y  & miss   & write X back, read Y      & Y, D$=$0 \\
    8    & read X  & miss   & read X                    & X, D$=$0 \\
    \bottomrule
  \end{tabular}
\end{table}

\section{Putting it together}
\label{sec:l12-summary}

A cache is fully described by a few design choices, from which its structure
follows.\source{Lesson 12, videos 41--45; H\&P §B.1.}
\begin{description}
  \item[Cache size and block size] give the number of lines and the offset
    bits.
  \item[Associativity] gives the number of sets, hence the index and tag bits
    (\cref{eq:l12-fields}).
  \item[Replacement policy] LRU adds a $\log_2 N$-bit counter to every line
    (none for direct-mapped caches).
  \item[Write policy] write-back adds a dirty bit to every line.
\end{description}
Every line thus stores a valid bit, a dirty bit, an LRU counter and a tag in
addition to its data.  \Cref{fig:l12-access-flow} shows how these are used on
one access.  A hit reads or writes the data and updates the LRU counters.  A
miss first chooses a victim, writes it back if it is valid and dirty, and
loads the new block; the access then proceeds as a hit.

\begin{figure}[htbp]
  \centering
  % One memory access in a set-associative cache with LRU replacement and
% the write-back policy.
\begin{tikzpicture}[x=1cm,y=1cm, font=\footnotesize\sffamily, >={Stealth[length=1.8mm]},
  stp/.style={draw, fill=ln-box, align=center, inner sep=3pt, minimum height=0.8cm, minimum width=4.4cm},
  miss/.style={stp, fill=white, minimum width=4.2cm},
  dec/.style={draw, diamond, aspect=2.4, fill=ln-accent!15, align=center, inner sep=1pt},
  lab/.style={font=\sffamily\scriptsize, text=ln-muted}]
  % --- hit path (left column)
  \node[stp] (split) at (0,0)
    {\iflangja{アドレスをタグ，インデックス，\\オフセットに分ける}{split the address into\\tag, index and offset}};
  \node[stp] (cmp) at (0,-1.3)
    {\iflangja{インデックスでセットを選び，\\有効な各ラインのタグと比較}{select the set by the index; compare\\the tag with every valid line of the set}};
  \node[dec] (hit) at (0,-2.65) {\iflangja{ヒット?}{hit?}};
  \node[stp] (acc) at (0,-6.55)
    {\iflangja{オフセットの位置を読み書き\\（書き込みなら D $\gets 1$）\\LRU カウンタを更新}{read or write at the offset\\(on a write: D $\gets 1$);\\update the LRU counters}};
  % --- miss path (right column)
  \node[miss] (victim) at (5.8,-2.65)
    {\iflangja{追い出すライン：無効なライン，\\なければ LRU カウンタが 0 のライン}{victim: an invalid line, or else\\the line whose LRU counter is 0}};
  \node[dec] (dirty) at (5.8,-3.95) {\iflangja{有効かつ\\ダーティ?}{valid and\\dirty?}};
  \node[miss] (wb) at (5.8,-5.25)
    {\iflangja{追い出すブロックを\\主記憶に書き戻す}{write the victim block\\back to main memory}};
  \node[miss] (load) at (5.8,-6.55)
    {\iflangja{ブロックを読み込みタグを格納\\V $\gets 1$，D $\gets 0$}{load the block, store its tag,\\V $\gets 1$, D $\gets 0$}};
  % --- arrows
  \draw[->] (split) -- (cmp);
  \draw[->] (cmp) -- (hit);
  \draw[->] (hit) -- node[lab, above] {\iflangja{いいえ}{no}} (victim);
  \draw[->] (hit) -- node[lab, right, pos=0.3] {\iflangja{はい}{yes}} (acc);
  \draw[->] (victim) -- (dirty);
  \draw[->] (dirty) -- node[lab, right] {\iflangja{はい}{yes}} (wb);
  \draw[->] (wb) -- (load);
  \draw[->] (dirty.east) -- node[lab, above] {\iflangja{いいえ}{no}} (8.4,-3.95) |- (load.east);
  \draw[->] (load) -- (acc);
\end{tikzpicture}

  \caption{One memory access in a set-associative cache with LRU replacement
    and the write-back policy.}
  \label{fig:l12-access-flow}
\end{figure}

\begin{example}
  \label{ex:l12-summary-trace}
  Consider the 2-way set-associative cache of \cref{ex:l12-sa-conflict}
  (16-byte blocks, index bits 5--4, tag bits 31--6), now with LRU replacement
  and the write-back policy.  \Cref{tab:l12-summary-trace} follows its set~1 through
  five accesses; both lines of the set are valid throughout.\tip{Memory
  traffic under write-back: one read per miss plus one write per eviction of
  a dirty block; hits cost nothing.}  The first
  access hits and only updates the counters.  The second misses, and its
  victim is way~1, whose counter is~0; that block is dirty and is written
  back before the new block is loaded.  The fourth access is a write miss
  whose victim is clean, so nothing is written back, and the loaded block
  becomes dirty at once.  The last access evicts the block with tag
  \texttt{0x40}, which the third access has written, and writes it back.
\end{example}

\begin{table}[htbp]
  \centering\small
  \caption{Set~1 of a 2-way set-associative cache with LRU replacement and
    write-back (\cref{ex:l12-summary-trace}): tag, dirty bit and LRU counter
    of each way after each access.  All tags, including those of the memory
    operations, are hexadecimal.}
  \label{tab:l12-summary-trace}
  \setlength{\tabcolsep}{4pt}
  \begin{tabular}{@{}llccccccl@{}}
    \toprule
    & & \multicolumn{3}{c}{Way 0} & \multicolumn{3}{c}{Way 1} & \\
    \cmidrule(lr){3-5}\cmidrule(lr){6-8}
    Access & Result & Tag & D & LRU & Tag & D & LRU & Memory operations \\
    \midrule
    (initial)             &      & 40 & 0 & 0 & 43 & 1 & 1 & \\
    read \texttt{0x1018}  & hit  & 40 & 0 & 1 & 43 & 1 & 0 & --- \\
    read \texttt{0x1054}  & miss & 40 & 0 & 0 & 41 & 0 & 1 & write back 43, load 41 \\
    write \texttt{0x101C} & hit  & 40 & 1 & 1 & 41 & 0 & 0 & --- \\
    write \texttt{0x1098} & miss & 40 & 1 & 0 & 42 & 1 & 1 & load 42 \\
    read \texttt{0x1058}  & miss & 41 & 0 & 1 & 42 & 1 & 0 & write back 40, load 41 \\
    \bottomrule
  \end{tabular}
\end{table}

\begin{example}
  \label{ex:l12-overhead}
  \Cref{tab:l12-organizations} works out the organization of a
  64~KB cache with 128-byte blocks, 32-bit addresses,
  LRU replacement and write-back, for several associativities.  For the
  8-way cache, the address \texttt{0x0040A7C6} (block number
  \texttt{0x814F}) has index $\texttt{0x814F} \mathbin{\&} \texttt{0x3F} = 15$
  and tag $\texttt{0x814F} \gg 6 = \texttt{0x205}$.  Its 512 lines need
  $512 \times 24 = \num{12288}$ bits of metadata, i.e.\ 1536 bytes,
  about \SI{2.3}{\percent} of the 64~KB of data.\margin{Metadata is a reason
  not to shrink blocks: with 16-byte blocks the same 8-way cache would have
  4096 lines of 24 bits, 12~KB or \SI{19}{\percent} of its data.}
\end{example}

\begin{table}[htbp]
  \centering\small
  \caption{Address fields and metadata bits per line (valid $+$ dirty $+$
    LRU counter $+$ tag) of a 64~KB cache with 128-byte blocks and 32-bit
    addresses.}
  \label{tab:l12-organizations}
  \begin{tabular}{@{}lrcccc@{}}
    \toprule
    Organization & Sets & Offset & Index & Tag & Metadata per line \\
    \midrule
    Direct-mapped     & 512 & 7 & 9 & 16 & $1+1+0+16 = 18$ \\
    2-way             & 256 & 7 & 8 & 17 & $1+1+1+17 = 20$ \\
    8-way             & 64  & 7 & 6 & 19 & $1+1+3+19 = 24$ \\
    Fully associative & 1   & 7 & 0 & 25 & $1+1+9+25 = 36$ \\
    \bottomrule
  \end{tabular}
\end{table}

Lesson 13 applies the same ideas---blocks, tags and replacement---to the
translation of virtual addresses, and Lesson 14 turns to techniques that
lower the hit time, the miss rate and the miss penalty of the basic cache
described here.

\begin{keypoint}[Lesson summary]
  Programs access memory with temporal and spatial locality, so a small, fast
  cache of recently used blocks serves most accesses.  Its benefit is
  measured by $\text{AMAT} = \text{hit time} + \text{miss rate} \times
  \text{miss penalty}$ (\cref{eq:l12-amat}).  Memory is divided into
  power-of-two blocks placed into cache lines; an address splits into tag,
  index and offset, and a line hits when it is valid and its tag matches.
  Direct-mapped, set-associative and fully associative caches differ in how
  many lines a block may occupy, trading hit time against miss rate.  On a
  miss LRU evicts the least recently used block, and the write-back policy
  with a dirty bit writes a block to memory only when a modified block is
  evicted.
\end{keypoint}

\end{document}
